% Gemini theme
% https://github.com/anishathalye/gemini
%
% We try to keep this Overleaf template in sync with the canonical source on
% GitHub, but it's recommended that you obtain the template directly from
% GitHub to ensure that you are using the latest version.

\documentclass[final]{beamer}

% ====================
% Packages
% ====================

\usepackage{amsmath, amsthm, amssymb, amsfonts}
\usepackage[size=custom,width=120,height=72,scale=1.0]{beamerposter}
\usetheme{gemini}
\usecolortheme{gemini}
\usepackage{bigints}
\usepackage{booktabs}
\usepackage{epigraph}
\usepackage{float}
\usepackage[T1]{fontenc}
\usepackage{graphicx}
\graphicspath{ {./img/} }
\usepackage{hhline}
\usepackage{hyperref}
\hypersetup{
colorlinks=true,
linkcolor=cyan,    
urlcolor=magenta,
pdftitle={Impure Relationalism in Classical Cosmology},
}
\usepackage{lmodern}
\usepackage{mathrsfs}
\usepackage{mathtools}
\usepackage{pgfplots}
\pgfplotsset{compat=1.14}
\usepackage{soul}
\usepackage{tcolorbox}
\usepckage{tensor}
\usepackage{tikz}
\usepackage{tikz-cd}
\usepackage[normalem]{ulem}

\newenvironment{highlight}
{\begin{tcolorbox}[colback=lightblue, boxrule=0pt, frame empty]}
{\end{tcolorbox}}

% SYMBOLS

% Display
\newcommand{\cp}{\clearpage}
\newcommand{\pb}{\pagebreak}

\newcommand{\dis}{\displaystyle}

% First-order logic
\newcommand{\exi}{\exists \:}
\newcommand{\fa}{\forall \:}
\newcommand{\Ra}{\Rightarrow}

% Algebra
\newcommand{\ca}{\cancel}
\newcommand{\rsa}{\rightsquigarrow}

% Set theory
% (Special sets)
\newcommand{\es}{\emptyset}
\newcommand{\N}{\mb{N}}
\newcommand{\Z}{\mb{Z}}
\newcommand{\Q}{\mb{Q}}
\newcommand{\R}{\mb{R}}
\newcommand{\C}{\mb{C}}

% (Operators)
\newcommand{\bs}{\backslash}
\newcommand{\mb}{\mathbb}
\newcommand{\mc}{\mathcal}

% SYMBOL SCHEMAS

% Display
\newcommand{\blue}[1]{\textcolor{blue}{#1}}
\newcommand{\Blue}[1]{\boxed{\blue{#1}}}
\newcommand{\plue}[1]{\textcolor{blue}{\pmb{#1}}}

\newcommand{\qua}[1]{\blue{#1} \:}

% Algebra
\newcommand{\canc}[2]{\tensor[^{#2}]{\cancel{#1}}{}}

% Bracket-based clusters
\newcommand{\p}[1]{\left( #1 \right)}
\newcommand{\s}[1]{\left[ #1 \right]}
\newcommand{\B}[1]{\left\{ #1 \right\}}
\newcommand{\pmat}[1]{\begin{pmatrix}
#1
\end{pmatrix}}
\newcommand{\smat}[1]{\begin{bmatrix}
#1
\end{bmatrix}}
\newcommand{\bmat}[1]{\begin{Bmatrix}
#1
\end{Bmatrix}}

\newcommand{\ucorn}[1]{\overset{\left\ulcorner \enspace \hphantom{#1} \enspace \right\urcorner}{#1}}
\newcommand{\lcorn}[1]{\overset{\left\llcorner \enspace \hphantom{#1} \enspace \right\lrcorner}{#1}}

% Modifiers
\newcommand{\wh}[1]{\widehat{#1}}
\newcommand{\wt}[1]{\widetilde{#1}}

% Set theory

% (Construction)
\newcommand{\set}[2]{\left\{ \vphantom{#2} #1 \right\rvert \left. \hspace{-2.3pt} \vphantom{#1} #2 \right\}}
\newcommand{\bsl}[2]{\left. \left\{ #1 \right\} \middle\bs \left\{ #2 \right\} \right.}
\newcommand{\bslset}[4]{\left. \set{#1}{#2} \middle\bs \set{#3}{#4} \right.}
\newcommand{\tup}[2]{\left( #1 \right\rvert \left. \hspace{-2.3pt} \vphantom{#1} #2 \right)}

% (Operators)
\newcommand{\pow}[1]{\mc{P}\p{#1}}

\newcommand{\im}[2]{\text{im}_{#1}\p{#2}}
\newcommand{\preim}[2]{\text{preim}_{#1}\p{#2}}

\newcommand{\id}[1]{\text{id}_{#1}}

% Structures

% (Geometry)
\newcommand{\inner}[2]{\left\langle #1, #2 \right\rangle}
\newcommand{\norm}[1]{\left\lvert #1 \right\rvert}

% Group theory

% (Operators)
\newcommand{\gen}[1]{\left\langle #1 \right\rangle}
\newcommand{\ord}[1]{\text{ord}\p{#1}}
\newcommand{\ordg}[2]{\text{ord}_{#1}\p{#2}}

% (Special groups)
\newcommand{\D}[2]{\text{D}_{#1}\p{#2}}
\newcommand{\GL}[2]{\text{GL}_{#1}\p{#2}}
\newcommand{\Or}[2]{\text{Or}_{#1}\p{#2}}
\newcommand{\SL}[2]{\text{SL}_{#1}\p{#2}}
\newcommand{\SO}[2]{\text{SO}_{#1}\p{#2}}
\newcommand{\SU}[2]{\text{SU}_{#1}\p{#2}}

% Category theory

% (Special categories)
\newcommand{\Cat}{\text{Cat}}
\newcommand{\Field}{\text{Field}}
\newcommand{\Grp}{\text{Grp}}
\newcommand{\Grph}{\text{Grph}}
\newcommand{\Met}{\text{Met}}
\newcommand{\Meas}{\text{Meas}}
\newcommand{\Mon}{\text{Mon}}
\newcommand{\Ord}{\text{Ord}}
\newcommand{\Ring}{\text{Ring}}
\newcommand{\Set}{\text{Set}}
\newcommand{\Top}{\text{Top}}

\newcommand{\Man}[1]{\text{Man}^{#1}}
\newcommand{\Mod}[1]{{#1}-\text{Mod}}
\newcommand{\Vect}[1]{\text{Vekt}_{#1}}

% (Operations)
\newcommand{\mor}[1]{\text{mor}\p{#1}}
\newcommand{\ob}[1]{\text{ob}\p{#1}}

% (Special morphisms)
% [On categories]
\newcommand{\aut}[1]{\text{aut}\p{#1}}
\newcommand{\epi}[1]{\text{epi}\p{#1}}
\newcommand{\inaut}[1]{\text{inn}\p{#1}}
\newcommand{\iso}[1]{\text{iso}\p{#1}}
\newcommand{\mono}[1]{\text{mono}\p{#1}}

% [On object pairs]
\newcommand{\Aut}[2]{\text{Aut}\p{#1, #2}}
\newcommand{\Epi}[2]{\text{Epi}\p{#1, #2}}
\newcommand{\Hom}[2]{\text{Hom}\p{#1, #2}}
\newcommand{\Inaut}[2]{\text{Inn}\p{#1, #2}}
\newcommand{\Iso}[2]{\text{Iso}\p{#1, #2}}
\newcommand{\Mono}[2]{\text{Mono}\p{#1, #2}}

% ====================
% Lengths
% ====================

% If you have N columns, choose \sepwidth and \colwidth such that
% (N+1)*\sepwidth + N*\colwidth = \paperwidth
\newlength{\sepwidth}
\newlength{\colwidth}
\setlength{\sepwidth}{0.025\paperwidth}
\setlength{\colwidth}{0.3\paperwidth}

\setlength\epigraphwidth{\colwidth}

\newcommand{\separatorcolumn}{\begin{column}{\sepwidth}\end{column}}

\newcommand{\bs}{\backslash}
\newcommand{\ca}{\cancel}
\newcommand{\dis}{\displaystyle}
\newcommand{\ex}{\exists \:}
\newcommand{\fa}{\forall \:}
\newcommand{\im}{\text{im}}
\newcommand{\mb}{\mathbb}
\newcommand{\mc}{\mathcal}
\newcommand{\mf}{\mathfrak}
\newcommand{\preim}{\text{preim}}
\newcommand{\ra}{\rightarrow}
\newcommand{\sing}{\text{Sing}}

\newcommand{\inc}[1]{\xhookrightarrow{#1}}
\newcommand{\mor}[1]{\xrightarrow{#1}}
\newcommand{\lmor}[1]{\xrightarrow[#1]}

\theoremstyle{remark}
\newtheorem*{remark}{Remark}

% ====================
% Title
% ====================

\title{Impure Relationalism in Classical Cosmology: Towards Energetic Causal Sets}

\author{Siddhartha Bhattacharjee}

\institute[shortinst]{University of Waterloo}

% ====================
% Footer (optional)
% ====================

\footercontent{
Blog (Tempus Spatium): \href{https://booodaness.github.io/tempus-spatium/}{https://booodaness.github.io/tempus-spatium/} \hfill
\href{https://cupc.cap.ca/}{Canadian Undergraduate Physics Conference} 2024, University of British Columbia \hfill
Email: \href{mailto:s3bhatta@uwaterloo.ca}{s3bhatta@uwaterloo.ca}}
% (can be left out to remove footer)

% ====================
% Logo (optional)
% ====================

% use this to include logos on the left and/or right side of the header:
% \logoright{\includegraphics[height=7cm]{logo1.pdf}}
% \logoleft{\includegraphics[height=7cm]{logo2.pdf}}

% ====================
% Body
% ====================

\begin{document}

\begin{frame}[t]
\begin{columns}[t]
\separatorcolumn

\begin{column}{\colwidth}

\epigraph{\textit{Nature exists only once. It is only our schematizing reproduction that creates equal cases. It is only in those that there exists mutual dependence of certain qualities.}}{Ernst Mach}

\begin{block}{Containment axioms}
\begin{alertblock}{Property containment}
Suppose $X, Y$ are in $\Set$. A property $Q$ on $X \times Y$ can be a \textbf{containment} of a property $P$ on $X$ in the following sense,

$$\qua{\fa P \in 2^X \: \fa Q \in 2^{X \times Y} :} \: Q \rhd P : \equiv \exi y \in Y : \fa x \in X : Q(x, y) \iff P(x)$$

We can recursively, using Cartesian products, extend the above definition to that of a $(X)^N = (X_1 \times \cdots \times X_N)$-property containing a $(X)^{N-k} = (X_1 \times \cdots \times X_{N-k})$-property,

\begin{align*}
\qua{\fa P \in 2^{(X)^{N-k}} \: \fa Q \in 2^{(X)^N} :} \: Q \rhd P & : \equiv \exi (x')_{N-k+1}^{N} \in (X)_{N-k+1}^{N} : \fa (x)^{N-k} \in (X)^{N-k} \\
& : Q((x)^{N-k} * (x')_{N-k+1}^N) \iff P((x)^{N-k})
\end{align*}

where $(x)_k^{k+l}$ denotes elements of $(X)_k^{k+l} = X_k \times \cdots \times X_{k+l}$ and $*$ is tuple concatenation.
\end{alertblock}

Informally, an $N$-parameter predicate $Q$ contains an $(N-k)$-parameter predicate $P$ iff it is possible to set the last $k$ parameters of $Q$ such that the new $(N-k)$-parameter predicate obtained via the free parameters of $Q$, is identical to $P$. \emph{If the said $k$ parameters can be arbitrarily set, $Q$ is said to be a \textbf{trivial containment} of $P$},

\begin{alertblock}{Trivial containment}
\begin{align*}
\qua{\fa P \in 2^{(X)^{N-k}} \: \fa Q \in 2^{(X)^N} :} & Q \unrhd P : \equiv \fa R \in [Q] : R \rhd P \text{ where,} \\
\qua{\fa Q, R \in 2^{(X)^N} :} & Q \sim R : \equiv \exi \underset{\text{(proj)}}{\pi} \in \p{(X)^N \to (X)^{N-k}} \\
& : \fa U, V \subseteq (X)^N : \p{\pi(U) = \pi(V)} \implies \p{Q \big\vert_U = R \big\vert_V}
\end{align*}
\end{alertblock}

Intuitively, the mechanism of the above axiom is that $\pi$ forgets the last $k$ parameters of $Q$, and if this does not change its extension which is modelled by $P$, then $Q$ trivially contains $P$. This allows us to \emph{'compress' any given $Q$ to some maximal \textbf{reduction} $\kappa(Q)$ with the same extensional information},

\begin{alertblock}{Property reduction}
\begin{align*}
\qua{\fa Q \in 2^{(X)^N} : \fa P \in \bigcup\limits_{k=0}^{N} 2^{(X)^k} :} & P \unrhd \kappa(Q) \\
\qua{\fa Q \in 2^{(X)^N} : \fa d \in \N :} & \p{\dim(Q) = d} : \iff \exi (Y)^d \subseteq (X)^N : \p{\kappa(Q) \in 2^{(Y)^d}}
\end{align*}
\end{alertblock}
\end{block}

where $(Y)^d \subseteq (X)^N$ denotes ordered string inclusion. Remarkably, \blue{$\p{\frac{\bigcup\limits_{k=0}^{N} 2^{(X)^k}}{\sim}, \succeq}$ \textbf{is a poset}}, where $\succeq$ is the operation induced by quotienting by $\unrhd$. 

\newline

By Zorn's lemma, \plue{every property has a unique reduction and dimension}.

\end{column}

\separatorcolumn

\begin{column}{\colwidth}

\begin{block}{Intrinsic and relational properties}
A property $Q \in 2^{(X)^N}$ is \textbf{$x_k$-intrinsic} for some $x_k \in X_k$ iff $Q(\bullet, \dots, \bullet, x_k, \bullet, \dots, \bullet)$ $\kappa$-reduces to some $P \in 2^{X_k}$,

\begin{alertblock}{Intrinsic properties as locally reducible}
\begin{align*}
\qua{\fa Q \in 2^{(X)^N}} & \qua{: \fa k \leq N : \fa x_k \in X_k} : Q \in \mc{I}_k(x_k) \\ 
: \iff \exi P \in 2^{X_k} & : P = \kappa(Q(\bullet, \dots, \bullet, x_k, \bullet, \dots, \bullet))
\end{align*}
\end{alertblock}

On the other hand, $Q$ is relational with respect to $(x_{k_1}, \dots, x_{k_d}) = (y)^d \in (Y)^d \subseteq (X)^N$ when,

\begin{alertblock}{Locally relational properties}
\begin{align*}
\qua{\fa Q \in 2^{(X)^N}} & \qua{: \fa d \in \N : \fa (Y)^d \subseteq (X)^N : \fa (x)^{(k)^d} \in (Y)^d :} Q \in \mc{R}_{k_1 \dots k_d}(x_{k_1}, \dots, x_{k_d}) \\ 
& : \iff \exi P \in 2^{(Y)^d} : P = \kappa(Q(\bullet, \dots, \bullet, x_{k_1}, \bullet, \dots, \bullet, x_{k_2}, \bullet, \dots, \bullet, x_{k_d}, \bullet, \dots, \bullet))
\end{align*}
\end{alertblock}

\blue{Due to the partial ordering properties of $\kappa$-reduction under appropriate bound variable quotienting, the rightmost existential clauses in the above definitions also turn out to involve \emph{uniqueness}}.
\end{block}

\begin{block}{Spatiotemporal locality}
In an event ontology, let $(X, \mc{O})$ be a topological space of events equipped with a \pmb{temporal partial ordering} $\preceq \in 2^{X \times X}$ and \pmb{causal structure} $\to \in 2^{X \times X}$. \textbf{Spatial locality} is the statement that \emph{for every pair of events with $x \to z$, there exists a causally relevant step $y$ for every open neighbourhood containing $x, y$},

\begin{highlight}
\begin{align*}
\qua{\fa x, z \in X :} S(x, z) & : \equiv \fa U \in \mc{O} : x, z \in U \implies \exi y \in U : x \to y \to z \\
S_X & : \equiv \qua{\fa x, z \in X :} x \to z \implies S(x, z)
\end{align*}
\end{highlight}

Similarly, \textbf{temporal locality} is the requirement that \emph{for every pair of causally linked events $x \to z$, it must be that $x \preceq z$ and the interval $[x, z] = \set{y \in X}{x \preceq y \preceq z}$ forms a \textbf{causal chain}},

\begin{highlight}
\begin{align*}
\qua{\fa x, z \in X :} T(x, z) & : \equiv (x \preceq z) \land \p{\fa y \in [x, z] : x \to y \to z} \\
T_X & : \equiv \qua{\fa x, z \in X :} (x \to z) \implies T(x, z)
\end{align*}
\end{highlight}

Lastly, \textbf{spatiotemporal locality} is the \emph{simultaneous requirement of spatial and temporal locality for each pair of points with either property},

\begin{highlight}
$$H_X : \equiv \qua{\fa x, z \in X :} (x \to z) \implies \s{S(x, z) \iff T(x, z)}$$
\end{highlight}

It follows from $T_X$ that $(X, \to)$ is a poset. If we prevent self-causation, it follows that $X$ is necessarily infinite. To prevent this restriction, one can upgrade $(X, \to)$ to a \textbf{causal set}, with the additional property of \textbf{local finiteness} $L_X$ i.e., \emph{every chain $[x, z]_{\to}$ is finite}.

The structure of a causal set $X$ with spatiotemporal locality is sufficiently characterized by the following,

\begin{alertblock}{Causation as primitive in causal sets}
$$L_X \land H_X \implies \begin{cases}
\mc{O} & = \set{[x, z]_{\to}}{x, z \in X} \\
\qua{\fa x, z \in X :} x \preceq z : & \equiv [x, z]_{\to} \neq \es 
\end{cases}$$
\end{alertblock}
\end{block}

\end{column}

\separatorcolumn

\begin{column}{\colwidth}

\begin{block}{Locality in second-order causation}
On top of causal set structure, one can have \textbf{property causation}, i.e., mechanisms of the form,

$$P(\bullet, \dots, \bullet, x_k, \bullet, \dots, \bullet) \to Q(\bullet, \dots, \bullet, z_l, \bullet, \dots, \bullet)$$

with the parameters being from topological spaces $(X_1, \mc{O}_1), \dots, (X_N, \mc{O}_N)$. To prevent trivial locality in the form of LHS and RHS being a local matter of fact at $z_l$, property locality must restrict $P$ to be $x_k$-intrinsic and $Q$ to be $z_l$-intrinsic, respectively,

\begin{alertblock}{Intrinsic residue in property causation}
The locality of $P(\dots, x_k, \dots) \to Q(\dots, z_l, \dots)$ is truth-evaluable iff $P \in \mc{I}_k(x_k)$ and $Q \in \mc{I}_l(z_l)$.
\end{alertblock}

Let us now see how \blue{classical dynamics hints at energy and momentum being the primitive intrinsic properties of a second-order causal set}.
\end{block}

\begin{block}{Relationalism in dynamical pairings}
Let $(X, \omega)$ be a symplectic manifold yielding the usual Poisson bracket $\{ \cdot, \cdot \}$. Given a Lagrangian $L$, position and conjugate momentum are \textbf{dynamically paired} on-shell, $\{ x^i, p_j \} = \delta\indices{^i_j}$. Noether's theorem leads to \plue{momenta being generators of translations}: if $x^i \to x^i + \delta x^i$, then, for a quantity $A(x^i, p_j)$, if $\delta A = \{ A, \delta G \}$, then $\exi \epsilon \in \R^+ : \delta G = \epsilon p_i$ i.e. $\delta A = \epsilon \{ A, p_i \}$. Similarly, \plue{energy is the generator of time-evolution}: $\delta G = \epsilon H$ implies that $\delta A = \epsilon \{ A, H \}$.

Consequently, since $\delta x^i$ is $X$-relational, $p_j$ is $P$-relational. \plue{But this only holds at the subsystem level}. In a cosmological setting, \emph{there can be no $X$-symmetries as any $x, y \in X$ are distinct iff they differ by intrinsic and relational properties},

\begin{alertblock}{Principle of identity of indiscernibles}
\begin{align*}
\qua{\fa k \leq N : \fa x_k, y_k \in X_k :} (x_k = y_k) & : \equiv \fa P \in \mc{R}_{\dots, k, \dots}(\dots, x_k, \dots) \cap \mc{R}_{\dots, k, \dots}(\dots, y_k, \dots) \\
& : P(\dots, x_k, \dots) \iff P(\dots, y_k, \dots)
\end{align*}
\end{alertblock}

Therefore, to model global dynamics on a causal set, \textbf{momentum and energy must be locally intrinsic properties}! This motivates the energetic causal set program, in which energy and momentum live as intrinsic quantities with causal efficacy on the nodes of $X$, and $\to$ dictates dynamics.
\end{block}

\begin{block}{References}
\nocite{*}
\footnotesize{\bibliographystyle{plain}\bibliography{poster}}
\end{block}

\end{column}

\end{columns}
\end{frame}

\end{document}